%
% Chapter 9: Parameter and Computation Analysis

\clearpage
\cleardoublepage

\chapter{Parameter and Computation Analysis}

\section{Importance of Model Efficiency}

Understanding model efficiency is crucial for:
\begin{itemize}
    \item Deployment on resource-constrained devices
    \item Training cost estimation
    \item Model selection and architecture design
    \item Benchmarking and comparison
\end{itemize}

Key metrics:
\begin{itemize}
    \item \textbf{Parameters:} Model size and memory footprint
    \item \textbf{FLOPs:} Computational complexity
    \item \textbf{Latency:} Inference time
    \item \textbf{Throughput:} Samples per second
\end{itemize}

\section{Parameter Counting}

Parameters are the learnable weights and biases in the network.

\subsection{Fully Connected Layer}

For input dimension $d_{in}$ and output dimension $d_{out}$:
\begin{align}
\text{Weights} &= d_{in} \times d_{out} \\
\text{Biases} &= d_{out} \\
\text{Total} &= d_{in} \times d_{out} + d_{out}
\end{align}

Example: FC layer from 784 to 256:
\begin{equation}
784 \times 256 + 256 = 200,704 + 256 = 200,960 \text{ parameters}
\end{equation}

\subsection{Convolutional Layer}

For input channels $C_{in}$, output channels $C_{out}$, kernel size $K$:
\begin{align}
\text{Weights} &= K \times K \times C_{in} \times C_{out} \\
\text{Biases} &= C_{out} \quad \text{(typically)} \\
\text{Total} &= K^2 \times C_{in} \times C_{out} + C_{out}
\end{align}

Example: Conv layer $3\times3$, 64 in, 128 out:
\begin{equation}
3 \times 3 \times 64 \times 128 + 128 = 73,728 + 128 = 73,856 \text{ parameters}
\end{equation}

\subsection{Batch Normalization}

For $C$ channels:
\begin{align}
\gamma \text{ (scale)} &= C \\
\beta \text{ (shift)} &= C \\
\mu \text{ (running mean)} &= C \quad \text{(inference only)} \\
\sigma \text{ (running var)} &= C \quad \text{(inference only)} \\
\text{Total trainable} &= 2C
\end{align}

Note: During training, BN has 2 trainable and 2 non-trainable parameters per channel.

\section{FLOPs (Floating Point Operations)}

FLOPs measure the computational cost of a forward pass.

\subsection{FLOPs for FC Layer}

For matrix multiplication $\mathbf{Y} = \mathbf{W}\mathbf{X} + \mathbf{b}$:
\begin{align}
\text{Multiplications} &= d_{out} \times d_{in} \quad \text{(per sample)} \\
\text{Additions} &= d_{out} \times (d_{in} - 1) \quad \text{(approximately } d_{out} \times d_{in}) \\
\text{FLOPs} &\approx 2 \times d_{out} \times d_{in} \quad \text{(counting mult+add as 2 FLOPs)}
\end{align}

For batch of $N$: FLOPs $= 2 \times N \times d_{out} \times d_{in}$

\subsection{FLOPs for Conv Layer}

For output feature map $H_{out} \times W_{out}$:
\begin{align}
\text{FLOPs per output pixel} &= K^2 \times C_{in} \times C_{out} \quad \text{(mult-add pairs)} \times 2 \\
\text{Total FLOPs} &= H_{out} \times W_{out} \times K^2 \times C_{in} \times C_{out} \times 2
\end{align}

Example: $224\times224$ input, $3\times3$ conv, 64 in, 128 out:
\begin{align}
H_{out} \times W_{out} &= 224 \times 224 = 50,176 \\
\text{FLOPs} &= 50,176 \times 3 \times 3 \times 64 \times 128 \times 2 \approx 1.17 \text{ GFLOPs}
\end{align}

\begin{note}[MACs vs FLOPs]
Sometimes FLOPs are reported as Multiply-Accumulate (MAC) operations:
\begin{equation}
\text{MAC} = \sum_{i} a_i \cdot b_i
\end{equation}
Each MAC typically counts as 1 operation (or 2 FLOPs if counting mult and add separately).

1 MAC $\approx$ 2 FLOPs for basic arithmetic.
\end{note}

\section{Practical Examples}

\subsection{LeNet-5 Parameter Count}

\begin{table}[h]
    \centering
    \begin{tabular}{L{2.5cm} L{2cm} L{2cm} L{2cm}}
        \toprule
        \textbf{Layer} & \textbf{Input} & \textbf{Weights} & \textbf{FLOPs} \\
        \midrule
        Conv1 & $1\times32\times32$ & $150$ & $2.7M$ \\
        Conv2 & $6\times14\times14$ & $2,400$ & $4.2M$ \\
        FC1 & 400 & $240K$ & $480K$ \\
        FC2 & 120 & $48K$ & $48K$ \\
        \midrule
        Total & - & $\sim$60K & $\sim$7M \\
        \bottomrule
    \end{tabular}
\end{table}

\subsection{AlexNet Parameter Count}

\begin{table}[h]
    \centering
    \begin{tabular}{L{2cm} L{2.5cm} L{2cm} L{2cm}}
        \toprule
        \textbf{Layer} & \textbf{Output} & \textbf{Params} & \textbf{FLOPs} \\
        \midrule
        Conv1 & $96\times55\times55$ & 35K & 105M \\
        Conv2 & $256\times27\times27$ & 614K & 448M \\
        Conv3 & $384\times13\times13$ & 885K & 149M \\
        Conv4 & $384\times13\times13$ & 1.3M & 224M \\
        Conv5 & $256\times13\times13$ & 884K & 149M \\
        FC6 & 4096 & 37.7M & 37.7M \\
        FC7 & 4096 & 16.8M & 16.8M \\
        FC8 & 1000 & 4.1M & 4.1M \\
        \midrule
        Total & - & $\sim$60M & $\sim$1.1B \\
        \bottomrule
    \end{tabular}
\end{table}

\section{Transfer Learning and Efficiency}

Transfer learning can dramatically reduce required data and training.

\begin{properties}[Transfer Learning Benefits]
\begin{itemize}
    \item Lower training cost (no training from scratch)
    \item Can use smaller datasets
    \item Faster convergence
    \item Often better final performance
\end{itemize}
\end{properties}

Common strategies:
\begin{enumerate}
    \item \textbf{Freeze:} Keep pretrained weights fixed
    \item \textbf{Fine-tune:} Adapt all/some layers with small learning rate
    \item \textbf{Extract features:} Use pretrained as feature extractor
\end{enumerate}

\section{Chapter Summary}

\begin{enumerate}
    \item Parameter count determines model size and memory requirements
    \item FLOPs measure computational complexity
    \item Efficient architectures like MobileNet reduce both
    \item Transfer learning leverages pretrained models efficiently
    \item Understanding efficiency guides model selection
\end{enumerate}
